\chapter*{Zusammenfassung}
\label{ch:zusammenfassung}

Neuronale Regelungsstrategien haben sich als wirksamer Ansatz zur Regelung
komplexer dynamischer Systeme etabliert und ermöglichen die Nachahmung modellprädiktiver Regler.
Dennoch ist ihr praktischer Einsatz in sicherheitskritischen Anwendungen nach wie vor durch das
Fehlen strenger Stabilitätsgarantien im geschlossenen Regelkreis eingeschränkt.
Diese Arbeit stellt einen Rahmen zum Lernen und zur formalen Zertifizierung neuronaler Imitationsregler
mithilfe gelernter Lyapunov-Funktionen vor.
Die vorgeschlagene Methodik kombiniert überwachtes Imitationslernen anhand eines MPC-Experten mit einem durch 
Falsifikation geleiteten gemeinsamen Training der Steuerungsstrategie und eines neuronalen Lyapunov-Kandidaten. 
Dabei werden extrahierte adversarische Gegenbeispiele verwendet, um die Steuerungsstrategie und den Lyapunov-Kandidaten 
hin zu einer verifikatorfreundlichen Stabilitätslandschaft zu verfeinern. 
Anschließend wird mithilfe des GPU-beschleunigten Verifikators $\alpha,\beta$-CROWN lokaler Stabilität und 
Vorwärtsinvarianz gesucht, um die Attraktionsbereich formal zu verifizieren. 
Die Pipeline wird sowohl an einem Doppelintegrator als auch an einem nichtlinearen Cart-Pole-System evaluiert.